\section{Transverse index derivatives and Gamma--Stieltjes gaps}
\label{tr:section}

A formula proved only after setting an index to $-m$ cannot be
differentiated in that index. The derivative sees information that a
Bernoulli polynomial does not contain. This section constructs that
information explicitly. We first obtain identities valid with an
independent moving index, and only then differentiate and specialize.
The kernel of Section~\ref{zbg:section} supplies a second, convergent
integral realization of the same derivatives.

\subsection{The joint Hurwitz completion}

Put
\begin{equation}\label{tr:Q}
 Q(u,x)=\zeta(u,x)-\frac1{u-1},
\end{equation}
with its removable value at $u=1$. This is entire in $u$ and holomorphic
for $\Re x>0$. At the resonant point,
\begin{equation}\label{tr:Qstieltjes}
 \left.\partial_u^rQ(u,x)\right|_{u=1}=(-1)^r\gamma_r(x).
\end{equation}
For $m,r\geq0$ set
\begin{equation}\label{tr:ell}
 \ell_{m,r}(x)=\left.\partial_u^r\zeta(u,x)\right|_{u=-m}.
\end{equation}
The derivatives of $Q$ at $-m$ differ from $\ell_{m,r}$ only by the
constant $r!/(m+1)^{r+1}$. Thus they give the same gap differences.

For $\Re s$ sufficiently large, uniformly on compact sets of $u$,
define the paired sum
\begin{equation}\label{tr:Dtilde}
 \widetilde D_{a,b}(s,u)=
 \sum_{n\geq0}\left\{
 (n+b)^{-s}Q(u,n+b)-(n+a)^{-s}Q(u,n+a)\right\}.
\end{equation}
We initially require convergence of the separate one-sided series,
not merely an accidental cancellation at an isolated exponent.

\begin{theorem}[A jointly entire weighted-tail difference]
\label{tr:jointD}
The continuation of \eqref{tr:Dtilde} is the entire function
\begin{equation}\label{tr:Djoint}
 \boxed{\widetilde D_{a,b}(s,u)
       =Q(u,b)h_{a,b}(s)-\HH_{a,b}(s,u).}
\end{equation}
Consequently the series
\begin{equation}\label{tr:D}
 D_{m,r}(s;a,b)=\sum_{n\geq0}\left\{
 (n+b)^{-s}\ell_{m,r}(n+b)
 -(n+a)^{-s}\ell_{m,r}(n+a)\right\},
 \qquad \Re s>m+2,
\end{equation}
has the entire continuation
\begin{equation}\label{tr:Didentity}
 D_{m,r}(s;a,b)=\ell_{m,r}(b)h_{a,b}(s)
       -\left.\partial_u^r\HH_{a,b}(s,u)\right|_{u=-m}.
\end{equation}
For each $r\geq0$, the similarly continued Stieltjes-weighted difference is
\begin{align}
 &\sum_{n\geq0}\left\{
 \frac{\gamma_r(n+b)}{(n+b)^s}
 -\frac{\gamma_r(n+a)}{(n+a)^s}\right\}
 \notag\\
 &\qquad=\gamma_r(b)h_{a,b}(s)
       -(-1)^r\left.\partial_u^r\HH_{a,b}(s,u)\right|_{u=1}.
 \label{tr:StieltjesD}
\end{align}
The last series is first interpreted, for example, on $\Re s>1$.
\end{theorem}

\begin{proof}
In a domain of absolute convergence, summing the inner index gives
\[
 Z_c(s,u)=\zeta(s,c)\zeta(u,c)
          -\sum_{n\geq0}(n+c)^{-s}\zeta(u,n+c).
\]
Insert this identity into
\[
 \HH_{a,b}(s,u)=Z_b(s,u)-Z_a(s,u)-\zeta(s,a)h_{a,b}(u).
\]
The endpoint products cancel to leave
\[
 \sum_{n\geq0}\left\{(n+b)^{-s}\zeta(u,n+b)
               -(n+a)^{-s}\zeta(u,n+a)\right\}
 =\zeta(u,b)h_{a,b}(s)-\HH_{a,b}(s,u).
\]
Subtract $h_{a,b}(s)/(u-1)$ to obtain \eqref{tr:Djoint}.
Its right side is jointly entire by Theorem~\ref{zbg:entire}.

The Hurwitz asymptotic expansion, uniformly differentiated on a compact
neighborhood of $u=-m$, gives
$\ell_{m,r}(n+c)=O(n^{m+1}(1+\log^r n))$ on compact endpoint sets.
A slightly enlarged compact set gives termwise spectral differentiation
in a sufficiently far right half-plane; continuation then proves
\eqref{tr:Didentity}. Formula \eqref{tr:Qstieltjes} gives
\eqref{tr:StieltjesD}. At $u=1$, the coefficient $\gamma_r(n+c)$ is
$O(1+\log^{r+1}n)$ on compact endpoint sets, proving the stated
initial half-plane.
\end{proof}

The whole function \eqref{tr:Djoint}, rather than two separately
regularized tails, fixes all constants. In particular the theorem covers
simultaneous spectral derivatives at $s=1,u=1$.

\subsection{A finite gap identity before specialization}

For lattice functions $f,g$ with enough decay, define
\begin{align}
 Z_c[f]&=\sum_{n\geq0}f(n+c),&
 Z_c[f,g]&=\sum_{n>j\geq0}f(n+c)g(j+c),\notag\\
 \HH_{a,b}[f,g]
   &=Z_b[f,g]-Z_a[f,g]-Z_a[f]\{Z_b[g]-Z_a[g]\}.
 \label{tr:decorated}
\end{align}
For longer decorated words, define
\[
 Z_c[f_1,\ldots,f_k]
   =\sum_{n_1>\cdots>n_k\geq0}\prod_{j=1}^k f_j(n_j+c),
\]
and define $\HH_{a,b}[f_1,\ldots,f_k]$ by the same
prefix--suffix recursion as \eqref{scope:relative}; the displayed
formula \eqref{tr:decorated} is its depth-two case.
All decorated expressions below mean these initial definitions and their
specified continuations. They do not mean arbitrary interpolations of
finite sums. Write $f_s(x)=x^{-s}$ and use pointwise products of letters.

\begin{theorem}[The moving-index gap and every transverse derivative]
\label{tr:gap}
The two functions
\[
 A(s,u,t)=\HH_{a,b}[f_sQ(u,\cdot),f_t],\qquad
 B(s,u,t)=\HH_{a,b}[f_s,f_tQ(u,\cdot+1)]
\]
continue individually to entire functions of $(s,u,t)$. They obey
\begin{align}
 A(s,u,t)
   &=\HH_{a,b}(u,s,t)+\HH_{a,b}(s+u,t)
                      +Q(u,a)\HH_{a,b}(s,t),\label{tr:A}\\
 B(s,u,t)&=A(s,u,t)+\HH_{a,b}(s,u,t).\label{tr:B}
\end{align}
Thus
\begin{equation}\label{tr:genericgap}
 \boxed{\HH_{a,b}(s,u,t)
   =\HH_{a,b}[f_s,f_tQ(u,\cdot+1)]
       -\HH_{a,b}[f_sQ(u,\cdot),f_t].}
\end{equation}
For $m,r\geq0$ this gives
\begin{align}
 J_{m,r}(s,t;a,b)
 &:=\left.\partial_u^r\HH_{a,b}(s,u,t)\right|_{u=-m}\notag\\
 &=\HH_{a,b}[f_s,f_t\ell_{m,r}(\cdot+1)]
      -\HH_{a,b}[f_s\ell_{m,r}(\cdot),f_t].
 \label{tr:J}
\end{align}
In particular this is an identity for arbitrary independent $s,t$ and
arbitrary right-half-plane endpoints.
\end{theorem}

\begin{proof}
First replace $Q$ by $\zeta$ and take all real parts large.
In the decorated sum $Z_c[f_s\zeta(u,\cdot),f_t]$, the additional
tail index is either larger than the outer index or equal to it.
Accordingly
\[
 Z_c[f_s\zeta(u,\cdot),f_t]
          =Z_c(u,s,t)+Z_c(s+u,t),
 \qquad
 Z_c[f_s\zeta(u,\cdot)]=Z_c(u,s)+\zeta(s+u,c).
\]
Substitute these identities in \eqref{tr:decorated}. Expanding the
relative recursion \eqref{scope:relative} and collecting the same cuts
gives
\[
 \HH_{a,b}[f_s\zeta(u,\cdot),f_t]
   =\HH_{a,b}(u,s,t)+\HH_{a,b}(s+u,t)
                            +\zeta(u,a)\HH_{a,b}(s,t).
\]
In the other decorated double, its two indices above the lower one
can be ordered in either direction or can be equal. These three cases
give $Z_c(s,u,t)$, $Z_c(u,s,t)$, and $Z_c(s+u,t)$.
The same relative recursion gives the extra term $\HH_{a,b}(s,u,t)$.

Subtracting the constant letter $1/(u-1)$ proves
\eqref{tr:A}--\eqref{tr:B}. Their right sides are entire, so they also
provide individual continuations of the decorated doubles through every
spectral intersection. Their difference proves \eqref{tr:genericgap}.
Differentiate this entire identity. At $u=-m$ the constants that
distinguish the derivatives of $Q$ from $\ell_{m,r}$ cancel between
the two doubles. This proves \eqref{tr:J}.
\end{proof}

The arbitrary-depth version of the counting argument will be useful
when several unmarked slots have first been removed. It also makes the
general-endpoint meaning of every boundary gap explicit.

\begin{lemma}[A decorated word with one moving tail]
\label{tr:masterlemma}
Let $\boldsymbol s=(s_1,\ldots,s_k)$ and $1\leq i\leq k$.
Then the decorated relative word has the jointly entire continuation
\begin{align}
 &\HH_{a,b}[f_{s_1},\ldots,f_{s_i}Q(u,\cdot),\ldots,f_{s_k}]
 \notag\\
 &\quad=Q(u,a)\HH_{a,b}(s_1,\ldots,s_k)
   +\sum_{j=0}^{i-1}\HH_{a,b}(s_1,\ldots,s_j,u,s_{j+1},\ldots,s_k)
 \notag\\
 &\qquad\quad
   +\sum_{j=1}^{i}\HH_{a,b}(s_1,\ldots,s_j+u,\ldots,s_k).
 \label{tr:masterdecorate}
\end{align}
In particular the top and bottom boundary identities are
\begin{align}
 \HH_{a,b}(u,\boldsymbol s)
  &=\HH_{a,b}[f_{s_1}Q(u,\cdot+1),f_{s_2},\ldots,f_{s_k}]
                      -Q(u,a)\HH_{a,b}(\boldsymbol s),
 \label{tr:topgap}\\
 \HH_{a,b}(\boldsymbol s,u)
  &=Q(u,b)\HH_{a,b}(\boldsymbol s)
      -\HH_{a,b}[f_{s_1},\ldots,f_{s_k}Q(u,\cdot)].
 \label{tr:bottomgap}
\end{align}
\end{lemma}

\begin{proof}
First replace $Q(u,x)$ by $\zeta(u,x)$ and take all real parts
large. In the infinite ordered sum, the extra tail index is at least
the index in position $i$. Its possible strict positions are the $i$
slots preceding that position, and its possible equalities are the
first $i$ existing indices. Thus the infinite decorated sum is exactly
the sum of the insertion terms and merger terms in
\eqref{tr:masterdecorate}, with $Z_c$ in place of $\HH$ and
without the endpoint term.

Expand each such $Z_b$ by deconcatenation and group terms by the cut
in the original undecorated word. For a cut before position $i$,
insertions and mergers lying wholly in the prefix give
$\zeta(u,a)Z_a(\text{prefix})$ by the ordinary singleton stuffle
identity. The remaining insertions and mergers lie in the suffix.
For a cut at or after position $i$, the same counting identity gives
$Z_a(\text{decorated prefix})$. These are exactly the two types
of cuts in the defining recursion for the decorated relative word.
Triangular uniqueness proves the claimed identity with $\zeta$;
subtracting $\HH(\boldsymbol s)/(u-1)$ replaces its endpoint value
by $Q(u,a)$ and proves \eqref{tr:masterdecorate}. Its right side
provides the jointly entire continuation.

At $i=1$, use $Q(u,x+1)=Q(u,x)-x^{-u}$ to obtain
\eqref{tr:topgap}. At $i=k$, the relative singleton stuffle identity
combines all insertions and mergers, and
$Q(u,b)-Q(u,a)=h_{a,b}(u)$ gives \eqref{tr:bottomgap}.
The difference of the decorations at adjacent positions $i+1$ and $i$,
using $Q(u,\cdot+1)$ at the lower position, leaves exactly the
insertion at the gap between them. This proves the interior gap
identity at any depth as well. No integer-endpoint uniqueness principle
is used.
\end{proof}

At an integer interval $a=b+N$, the meaning of the gap is especially
transparent:
\begin{equation}\label{tr:finitegap}
 J_{m,r}(s,t;b+N,b)=
 \sum_{N>i>k\geq0}
 \frac{\ell_{m,r}(k+b+1)-\ell_{m,r}(i+b)}
      {(i+b)^s(k+b)^t}.
\end{equation}
Indeed the middle-index sum is
\[
 \sum_{k<j<i}(j+b)^m(-\log(j+b))^r
           =\ell_{m,r}(k+b+1)-\ell_{m,r}(i+b).
\]
The proof above establishes the general-endpoint identity independently
of this finite illustration.

The Hurwitz shift identity also gives an equivalent form using only
unshifted decorated letters:
\begin{align}
 J_{m,r}(s,t;a,b)
   ={}&\HH_{a,b}[f_s,f_t\ell_{m,r}]
       -\HH_{a,b}[f_s\ell_{m,r},f_t]
       -\partial_t^r\HH_{a,b}(s,t-m).
 \label{tr:commutator}
\end{align}
The last term is essential. For $m=r=0$, it recovers exactly
\[
 \HH(s,0,t)=\HH(s-1,t)-\HH(s,t-1)-\HH(s,t),
\]
the strict zero-gap deletion rule.

\subsection{Gamma gaps, Stieltjes primitives, and repeated primitives}

For the first transverse derivative at zero, \eqref{scope:hurwitz}
turns \eqref{tr:finitegap} into
\begin{equation}\label{tr:logGammafinite}
 \left.\partial_u\HH_{b+N,b}(s,u,t)\right|_{u=0}
 =
 \sum_{N>i>k\geq0}
 \frac{\log\Gamma(k+b+1)-\log\Gamma(i+b)}
      {(i+b)^s(k+b)^t}.
\end{equation}
The constants $\log(2\pi)/2$ disappear. This is the requested transverse
identity in explicit Gamma coordinates, with \eqref{tr:J} specifying
its general-endpoint counterpart.

Differentiating $-\!u\zeta(1+u,x)$ at $u=0$ gives, for $r\geq1$,
\begin{equation}\label{tr:Stieltjesprimitive}
 \frac{d}{dx}\ell_{0,r}(x)=(-1)^r r\,\gamma_{r-1}(x).
\end{equation}
Therefore
\begin{equation}\label{tr:integratedgap}
 \ell_{0,r}(z+1)-\ell_{0,r}(y)
   =(-1)^{r+1}r\int_{z+1}^{y}\gamma_{r-1}(x)\dd x.
\end{equation}
The path lies in the right half-plane. The integral is path independent,
and its additive normalization is fixed by the displayed endpoints.
Combining \eqref{tr:J} with \eqref{tr:integratedgap} gives all the
corresponding harmonic identities with Stieltjes-primitive gaps.
At $u=1$, the same entire identity instead gives
\begin{align}
 \left.\partial_u^r\HH(s,u,t)\right|_{u=1}
 =(-1)^r\bigl\{\HH[f_s,f_t\gamma_r(\cdot+1)]
                 -\HH[f_s\gamma_r,f_t]\bigr\}.
 \label{tr:Stieltjesgap}
\end{align}
These statements retain the ordered information; they do not claim
a reduction of every resulting value to depth-one constants.

A standard balanced normalization packages repeated primitives of
log-Gamma \cite{EspinosaMoll}:
\begin{equation}\label{tr:balanced}
 \mathcal B_m(x)=
       \frac{\zeta'(-m,x)+H_m\zeta(-m,x)}{m!},\qquad H_0=0.
\end{equation}
Here $\mathcal B_0(x)=\log\Gamma(x)-\tfrac12\log(2\pi)$, and
\begin{equation}\label{tr:balancedderivative}
 \mathcal B_m'(x)=\mathcal B_{m-1}(x)\qquad(m\geq1).
\end{equation}
To check this, differentiate the Hurwitz identity
$\partial_x\zeta(u,x)=-u\zeta(u+1,x)$ once in $u$, then use
$H_m=H_{m-1}+1/m$. The factor $m!$ in \eqref{tr:balanced} is
exactly the one needed for \eqref{tr:balancedderivative}.

It follows that the repeated-Gamma-primitive gap has the exact identity
\begin{align}
 &\HH[f_s,f_t\mathcal B_m(\cdot+1)]
                 -\HH[f_s\mathcal B_m,f_t]\notag\\
 &\qquad=\frac{J_{m,1}(s,t)+H_m\HH(s,-m,t)}{m!}.
 \label{tr:balancedgap}
\end{align}
The second term is a finite Bernoulli reduction. Thus no undetermined
polynomial or integration constant is hidden in the primitive ladder.

As one elementary Stieltjes check, let $h=h_{a,b}$. Differentiating the
stuffle of a singleton with $\HH(s,t)$ gives
\begin{align}
 &\left.\partial_u\{\HH(u,s,t)+\HH(s,u,t)+\HH(s,t,u)\}\right|_{u=0}
 \notag\\
 &\qquad=h'(0)\HH(s,t)-(\partial_s+\partial_t)\HH(s,t).
 \label{tr:insertionderivative}
\end{align}
At $s=t=1$, use $\HH(v,v)=(h(v)^2-h(2v))/2$ to obtain the fully
specified expression
\begin{equation}\label{tr:Stieltjesexample}
 \frac{h'(0)}2\{h(1)^2-h(2)\}
   +h(1)\{\gamma_1(b)-\gamma_1(a)\}
   +\zeta'(2,b)-\zeta'(2,a).
\end{equation}
This is a consistency specialization of the classical stuffle algebra,
not a separate claim of a new product principle.

\subsection{Absolutely convergent transverse Mellin formulas}

The following elementary normalized-Mellin lemma makes the transverse
construction directly computable.

\begin{lemma}[Every derivative at a nonpositive Mellin order]
\label{tr:Mellinlemma}
Suppose $f$ is smooth on $[0,\infty)$, extends holomorphically near zero,
and all its derivatives decrease exponentially at infinity. Let
\[
 \mathcal M_f(u)=\frac1{\Gamma(u)}
                  \int_0^\infty y^{u-1}f(y)\dd y,\qquad\Re u>0.
\]
It has an entire continuation. Put $c_j=[v^j]\Gamma(v)^{-1}$.
For $m\geq0$,
\begin{align}
 \mathcal M_f(-m)&=(-1)^m f^{(m)}(0),\label{tr:Mvalue}\\
 \mathcal M_f^{(r)}(-m)
 &=(-1)^m r!\sum_{j=1}^r\frac{c_j}{(r-j)!}
      \int_0^\infty\frac{\log^{r-j}y}{y}
          \{f^{(m)}(y)-f^{(m)}(0)e^{-y}\}\dd y ,
 \qquad r\geq1.
 \label{tr:Mjet}
\end{align}
All integrals in \eqref{tr:Mjet} are ordinary absolutely convergent
integrals.
\end{lemma}

\begin{proof}
Repeated integration by parts gives
$\mathcal M_f(u)=(-1)^m\mathcal M_{f^{(m)}}(u+m)$, initially on
$\Re u>0$ and then by continuation. It is enough to expand at zero.
Write $g=f^{(m)}$ and use the exact identity
\[
 \mathcal M_g(v)=g(0)+\frac1{\Gamma(v)}
       \int_0^\infty y^{v-1}\{g(y)-g(0)e^{-y}\}\dd y .
\]
The last integral is holomorphic for $\Re v>-1$, locally at zero.
Its $q$th derivative at zero is the integral with
$\log^qy/y$. The bracket is $O(y)$ near zero and decreases
exponentially at infinity. Coefficient multiplication by
$\Gamma(v)^{-1}$ proves \eqref{tr:Mjet}.
The continuation and value formula also follow from repeated integration
by parts with an arbitrarily large number of derivatives.
\end{proof}

The coefficients $c_j$ are explicit Gamma Taylor coefficients:
\[
 c_1=1,\qquad c_2=\gamma,\qquad
 c_3=\frac{\gamma^2-\zeta(2)}2,
\]
and in every degree they follow from
\[
 \frac1{\Gamma(v)}
     =v\exp\left(\gamma v+
             \sum_{q\geq2}\frac{(-1)^{q+1}\zeta(q)}qv^q\right).
\]

\begin{theorem}[Convergent formula for every transverse gap]
\label{tr:integral}
Let $\Re s,\Re t>0$, $m\geq0$, and $r\geq1$. In
\eqref{tr:Mjet}, set
\[
 f(y)=\mathscr K_{a,b}^{(3)}(x,y,z).
\]
Then $J_{m,r}(s,t;a,b)$ is the normalized $(x,z)$ Mellin transform
of its right side. Explicitly it is
\begin{align}
 &\frac{(-1)^m r!}{\Gamma(s)\Gamma(t)}
 \sum_{j=1}^r\frac{c_j}{(r-j)!}
 \int_0^\infty\int_0^\infty
       x^{s-1}z^{t-1}\notag\\
 &\qquad\cdot\int_0^\infty \frac{\log^{r-j}y}{y}
 \left\{\partial_y^m\mathscr K^{(3)}_{a,b}(x,y,z)
       -\partial_y^m\mathscr K^{(3)}_{a,b}(x,0,z)e^{-y}\right\}
             \dd y\dd z\dd x .
 \label{tr:tripleintegral}
\end{align}
All these integrals are absolutely convergent. Their common continuation
is entire in $s,t$. Every further fixed spectral derivative can be
taken in the convergent integrals on compact subsets of
$\Re s,\Re t>0$.
\end{theorem}

\begin{proof}
Lemma~\ref{zbg:bounds} gives uniform exponential control of all derivatives
of the kernel, including at $y=0$. The inner bracket is uniformly
$O(y)e^{-\delta(x+z)}$ as $y\downarrow0$; for $y\geq1$ it is bounded
by $Ce^{-\delta(x+z)}(e^{-\delta y}+e^{-y})$.
Consequently the inner logarithmic integrals and the two outer Mellin
integrals are absolutely convergent. Fubini's theorem and the preceding
lemma give the derivative of
$\mathscr E_{a,b}(s,u,t;\boldsymbol0)=\HH_{a,b}(s,u,t)$.
Theorem~\ref{zbg:entire} gives the entire continuation. Powers of $\log x$
and $\log z$ from further spectral derivatives preserve the same
integrable majorants on compact half-plane subsets.
\end{proof}

For $m=0,r=1$ the formula contains just
\[
 \int_0^\infty
 \frac{\mathscr K^{(3)}(x,y,z)-\mathscr K^{(3)}(x,0,z)e^{-y}}y\dd y.
\]
This supplies a convergent realization of the log-Gamma gap at arbitrary
noninteger endpoints, without separate divergent ordered sums.

\subsection{All negative outer values and fixed Barnes constants}

The weighted difference in Theorem~\ref{tr:jointD} has especially
explicit values at nonpositive outer orders.

\begin{theorem}[Finite Bernoulli formula for the transverse weighted tails]
\label{tr:negativeD}
Let $d,m,r\geq0$ and $q=d+1$. Then
\begin{align}
 D_{m,r}(-d;a,b)
   =\frac1q\biggl\{&
        B_q(a)\ell_{m,r}(a)-B_q(b)\ell_{m,r}(b)\notag\\
       &+\sum_{k=0}^q\binom qk B_{q-k}(1)
                             h_{a,b}^{(r)}(-m-k)\biggr\}.
 \label{tr:negativeformula}
\end{align}
Every term is an endpoint Hurwitz derivative or a Bernoulli polynomial.
\end{theorem}

\begin{proof}
The required top-index deletion identity is
\begin{equation}\label{tr:topdelete}
 \HH(-d,u)=\frac1q\left\{
 B_q(a)h(u)-\sum_{k=0}^q\binom qk B_{q-k}(1)h(u-k)\right\}.
\end{equation}
For a direct analytic derivation, sum the leading index of $Z_c(v,u)$
as $\zeta(v,n+c+1)$, then continue $v$ to $-d$ while $\Re u$ is
large. Use
$\zeta(-d,x+1)=-B_q(x+1)/q$ and expand the Bernoulli polynomial.
Insert the result in the relative quotient. The leading singleton
$Z_a(-d)=-B_q(a)/q$ contributes the first term. This proves
\eqref{tr:topdelete} on a nonempty domain, and entireness in $u$ extends it.

Differentiate \eqref{tr:topdelete} $r$ times in $u$ and set $u=-m$.
Substitute in \eqref{tr:Didentity} and use
$h(-d)=(B_q(a)-B_q(b))/q$. The terms involving
$B_q(a)\ell_{m,r}(b)$ cancel, giving \eqref{tr:negativeformula}.
\end{proof}

For $m=0,r=1,d=0$, the standard Barnes relation
\cite{Vardi,DLMFBarnes} gives the particularly simple identity
\begin{equation}\label{tr:Barnes}
 \boxed{D_{0,1}(0;a,b)=
    \log G(a)-\log G(b)-\frac{a-b}{2}\log(2\pi).}
\end{equation}
Here $G$ is Barnes' function with its standard asymptotic normalization,
$G(1)=1$, and $G(x+1)=\Gamma(x)G(x)$. Holomorphic logarithms in the
right half-plane are normalized from the positive real axis.
Indeed the $q=1$ formula is
\[
 (a-1)\ell_{0,1}(a)-(b-1)\ell_{0,1}(b)
                    +\zeta'(-1,b)-\zeta'(-1,a),
\]
and
$\log G(x)=(x-1)\log\Gamma(x)+\zeta'(-1)-\zeta'(-1,x)$
proves \eqref{tr:Barnes}. At an integer endpoint difference it also
agrees with the finite product obtained by iterating the functional
equation. General $d$ gives the analogous finite endpoint formula in
higher Hurwitz derivatives, without choosing an unspecified multiple
Gamma normalization.

\subsection{A finite family after arbitrary unmarked deletions}

The preceding construction is stable under inserting any finite word
of unmarked nonpositive indices. We state exactly the scope.

\begin{proposition}[One marked index with arbitrary nonpositive decoration]
\label{tr:finitefamily}
Suppose a word contains $k\geq1$ independent free indices, one marked
index $u$, and unmarked indices $-m_1,\ldots,-m_e$, with $m_i\geq0$.
Put $M=\sum_i(m_i+1)$. After differentiating the marked index $r$ times
and setting $u=-m$, the word is a finite polynomial-endpoint combination
of $k$-letter relative sums whose letters are ordinary shifted powers,
ordinary spectral derivatives of those powers, and at most one factor
\[
 \ell_{m+\nu,r}(x),\qquad 0\leq\nu\leq M.
\]
Endpoint values of the same finite family may occur. Thus the new gap
functions needed for this one-marked problem are a specified finite
list. This is a bound for this representation, not an assertion of
minimal arithmetic depth.
\end{proposition}

\begin{proof}
First regard $u$ as free. The finite Bernoulli deletion rules for each
unmarked $-m_i$ hold identically in all remaining indices: sum its
finite strict gap by a Bernoulli polynomial, or use the leading-tail
argument of \eqref{tr:topdelete} at the upper boundary, and insert the
result in the relative recursion. An interior deletion has the two
endpoint terms
\[
 \frac1{m_i+1}\sum_{j=0}^{m_i+1}\binom{m_i+1}{j}
 \left\{B_{m_i+1-j}\HH(\ldots,s-j,t,\ldots)
       -B_{m_i+1-j}(1)\HH(\ldots,s,t-j,\ldots)\right\}.
\]
It shifts adjacent surviving orders down by nonnegative integers.
The total shift budget is at most $M$, even when adjacent unmarked
indices are combined during the deletion. Coefficients are endpoint
polynomials and are independent of every surviving free order.

Now differentiate in $u$ and set $u=-m$. Each surviving marked order
is $-m-\nu$ for some $0\leq\nu\leq M$.
Apply Lemma~\ref{tr:masterlemma} at the surviving marked position.
Its adjacent-decoration difference gives the interior gap, and
\eqref{tr:topgap}--\eqref{tr:bottomgap} give the boundary gaps.
Their finite-interval forms are
$\ell(y+1)-\ell(a)$ at the top and
$\ell(b)-\ell(y)$ at the bottom. Only one decorated factor is introduced.
Finally
$\ell_{m+\nu,r}(x+1)=\ell_{m+\nu,r}(x)
                   -x^{m+\nu}(-\log x)^r$
removes shifted potentials in favor of ordinary free-order derivatives.
The analytic identities at every step were proved with the marked
order still free, so the transverse differentiation is legitimate.
\end{proof}

There is also an analytic all-decoration version. Take any fixed
free word with one marked slot, apply Theorem~\ref{zbg:main}, then
differentiate that slot at $-m$. Local normality proves convergence on
the same full unit polydisc. Formula \eqref{tr:Mjet}, applied to the
marked kernel coordinate, gives the corresponding logarithmic-integral
generator. This simultaneously extends R7 and R8 to every derivative
order and every nonpositive marked center.

If there are no other free indices, the remaining singleton is
$h_{a,b}^{(r)}(-m-\nu)$, so no decorated lower-depth sum is needed.
Several independently marked indices are a further problem: this
proposition makes no claim that their repeated gap primitives close
in the same one-marked list.

